% GENERATED FILE. Edit canonical lessons, then run npm run generate:books.
% canonical-source-hash: fnv1a64:fb49102a7d9fbcc6
% canonical-lessons: FR-C22-chien-chat

\chapter{Dog and Cat}
\label{ch:fr-dog-and-cat}
\hlchaptermodality{22}

\begin{chapteropening}
\textbf{By the end of this chapter:} \emph{I can name a dog and a cat in French and say why chien is the regular descendant of canis while chat is a word that travelled out of Egypt.}

Say chien and check its ca- $\to$ cha- shift against champ and chanter, then say chat and place it beside Spanish gato and Italian gatto as one inherited cattus. The chapter's only lesson, so the payoff covers everything it introduces.
\end{chapteropening}

\section[chien, chat]{chien, chat — the regular one, and the well-traveled one}

\label{lesson:FR-C22-chien-chat}



\begin{quote}
You've met Spanish's \emph{perro}, a genuine unsolved mystery, and Latin's own \emph{canis}/\emph{fēlēs}$\to$\emph{cattus} history. French's two animal words tell a cleaner story: one is perfectly regular, the other is the same well- documented word you already know.
\end{quote}

\begin{cousinweb}
\textbf{chien} (\textquotedblleft{}\textbf{dog}\textquotedblright{}) comes straight from Latin \textbf{canis} (accusative \emph{canem}) — and unlike Spanish, French kept it. Two separate, regular changes are stacked here: the initial \textbf{c}$\to$\textbf{ch} is the same \textbf{ca- $\to$ cha-} palatalization you can check against other Latin words — \textbf{campus} $\to$ \textbf{champ} (\textquotedblleft{}field\textquotedblright{}), \textbf{cantāre} $\to$ \textbf{chanter} (\textquotedblleft{}to sing\textquotedblright{}); the \textbf{a}$\to$\textbf{ie} in the middle is a different, later shift specific to words like \emph{canem} (known as Bartsch's Law — the outcome is completely regular and undisputed, even though the precise phonetic trigger is still debated). Put together, \emph{chien} is a fully \textbf{regular} sound change, following French's own rules predictably — a genuinely different story from Spanish, where the expected descendant (\emph{can}) got pushed aside by the still-unexplained \emph{perro}.
\end{cousinweb}

\begin{cousinweb}
\textbf{chat} (\textquotedblleft{}\textbf{cat}\textquotedblright{}) continues \textbf{cattus}, the same Late Latin word — most likely ultimately \textbf{Afro-Asiatic} (compare Nubian \emph{kadis}) — that spread along Roman trade routes as domestic cats themselves spread out of \textbf{Egypt}. This is the identical word behind Spanish's \emph{gato} and Italian's \emph{gatto} — French, Spanish, and Italian all independently inherited the exact same Late Latin replacement for Classical \emph{fēlēs}.
\end{cousinweb}

\subsection*{Your turn}
Say these aloud:

\begin{itemize}
\raggedright
  \item \textquotedblleft{}chien\textquotedblright{} — dog, the regular descendant of canis
  \item \textquotedblleft{}champ, chanter\textquotedblright{} — the same ca- $\to$ cha- shift, in field and to sing
  \item \textquotedblleft{}chat\textquotedblright{} — cat, the same word as Spanish's gato, from Egypt via Rome
\end{itemize}

\begin{tcolorbox}[breakable,colback=teal!4,colframe=teal!35!black,title={Before you move on}]
Is \emph{chien} a regular or irregular descendant of Latin \emph{canis}? (\textbf{Regular} — the same \emph{ca- $\to$ cha-} shift as \emph{champ} and \emph{chanter}.) Does French show the same \emph{perro}-style mystery Spanish does for \textquotedblleft{}dog\textquotedblright{}? (\textbf{No} — French kept the expected Latin word.) What word does \emph{chat} continue, and where did that word most likely originate? (\textbf{Cattus}, most likely \textbf{Afro-Asiatic}, traveling with the cat out of \textbf{Egypt}.)
\end{tcolorbox}
